% ── §7d Methods: Comparison with Deep Learning Approaches ──
\section{Relationship to Deep Learning Approaches}
\label{sec:methods-dl}

The success of AlphaFold2~\cite{jumper2021} in protein structure prediction
has largely superseded standalone contact prediction as a practical tool.
However, we argue that the Karnaugh-map framework addresses a complementary
need that end-to-end deep learning does not: interpretability and formal
verifiability of the encoding layer.

AlphaFold2 processes MSAs through attention-based neural networks whose
internal representations are not decomposable into human-readable rules.  In
contrast, our Quine--McCluskey-minimised circuits produce sum-of-products
expressions such as ``contact if group$_i \in \{\text{hydrophobic},
	\text{aromatic}\}$ AND group$_j \in \{\text{hydrophobic}, \text{aromatic}\}$
AND separation $\in [6, 21)$'' that can be directly inspected, validated
against biophysical knowledge, and formally verified.

Furthermore, the formal verification layer (118 Lean~4 theorems) provides a
guarantee that no implementation bug corrupts the encoding, a property that no
neural network can offer.  While the absolute precision of our circuits
(0.179) is far below AlphaFold2's near-perfect predictions, the two approaches
address different questions: AlphaFold2 answers ``what is the structure?''
while our framework answers ``which specific residue-pair combinations are
coevolutionarily constrained, and can we prove that our method for detecting
them is correct?''

A promising future direction is to use circuit-derived features as input to
neural architectures, combining the interpretability of Boolean logic with
the representational power of deep learning.  The rank-fusion result
(Section~\ref{sec:results-fusion}) already demonstrates that K-map features
carry information orthogonal to continuous coupling scores; incorporating them
into a learned model might yield further improvements.
